\documentclass[11pt,a4paper]{article}

\errorcontextlines=9999

\usepackage[T1]{fontenc}
\usepackage[margin=2cm]{geometry}

\usepackage{libertinus}
\usepackage[scale=.95]{inconsolata}

\usepackage[prefix=]{xcolor-material}
\usepackage[colorlinks=true,allcolors=purple]{hyperref}

\usepackage{enumitem}
\usepackage{microtype}
\usepackage{parskip}

\usepackage{mathtools}
\usepackage{absint}

\begin{document}
Hey! This is: \absexpr{\lor}. \absexpr{\iff}.
Every symbol links to its entry in the legend at the end of this document.

\absexpr{a \lor b \asdef\iff c \eq 3 \asdef\eq 4}\\
\(a \lor b \iff c\)

Lattice: A lattice \absexpr{(X, \partof, \lub, \glb)} is a poset (partially ordered set), such that \absexpr{\forall a, b \in X: a \lub b} and \absexpr{a \glb b} exist (Mine).

Complete Lattice: A complete lattice \absexpr{(X, \partof, \lub, \glb, \bot, \top)} is a lattice, such that
\begin{enumerate}[nosep]
   \item \absexpr{\forall S \subseteq X: \Lub S} and \absexpr{\Glb S} exist.
   \item \absexpr{X} has a least element \absexpr{\bot} and a greatest element \absexpr{\top}.
\end{enumerate}

For any set \absexpr{X}, the powerset \absexpr{(\powerset(X), \subseteq, \union, \intersect, \emptyset, X)} is a complete lattice.

Lets write a fancy \absexpr{\partof_b} or \absexpr{R \glb_b X, \glb_b} or \absexpr{\lub_\#}!

\section{Concrete and abstract}
Single objects are marked with \verb|\abstract| and \verb|\concrete|:
\absexpr{\abstract{X} \partof \concrete{Y}}.
Whole expressions switch worlds with \verb|\abstractly| and \verb|\concretely|:
\absexpr{\concretely{(X, \partof, \lub, \glb, \bot, \top)}} becomes
\absexpr{\abstractly{(X, \partof, \lub, \glb, \bot, \top)}}, and
\absexpr{\abstractly{\Lub_{i \in \N} X_i \partof \widen}} keeps the mark on
the operators that have one.

A Galois connection \absexpr{\galoisbetween{\powerset(\Z)}{\NumericDomain}}
and a Galois insertion \absexpr{\powerset(\Z) \galoisinsertion \NumericDomain}
relate the two worlds via \absexpr{\abstraction} and \absexpr{\concretization},
with \absexpr{\abstraction \compose \concretization \partof \id}.

\section{Semantics}
\absexpr{\EvaluateExpr{x + 1}\Environment \eq \envlookup{\Environment}{x} \intadd 1}
for an environment \absexpr{\Environment \in \Environments \eq \Variables \to \NumericDomain},
and \absexpr{\abstractly{\EvaluateStatement{s}} \compose \abstraction \partof \abstraction \compose \EvaluateStatement{s}}.

The fixpoint of a statement is \absexpr{\lfp_{\bot} f \eq \Lub_{i \in \N} f^i(\bot)},
its abstract counterpart is \absexpr{\abstractly{\lfp_{\bot} f}} and it is reached
with a widening \absexpr{a \widen b} and a narrowing \absexpr{a \narrow b}.

\section{Sets}
\absexpr{\Set{x \in \N \Given x \leq 3} \union \Set{4}}, \absexpr{\card{X}},
\absexpr{\Tuple{a, b}}, \absexpr{\Interval{a}{b}}, \absexpr{\restrict{f}{S}},
\absexpr{\envupdate{\Environment}{x}{v}}.

\section{Lambda calculus}
\absexpr{(\lam{x}{M}) \app N \betared \subst{M}{x}{N}} and
\absexpr{\lam{x}{M \app x} \etared M} if \absexpr{x \notin \FV{M}}.
The fixpoint combinator unfolds as \absexpr{\combY \app f \betareds f \app (\combY \app f)},
and \absexpr{\combS \app \combK \app \combK \betaeq \combI}.
Church numerals: \absexpr{\church{3} \alphaeq \lam{f\,x}{f \app (f \app (f \app x))}}.

Typing: \absexpr{\typing{\ctxadd{\Context}{x}{\sigma}}{M}{\tau} \implies \typing{\Context}{\lam{x}{M}}{\sigma \to \tau}},
polymorphism: \absexpr{\tyabs{X}{\lamty{x}{X}{x}} \hastype \forallty{X}{X \to X}},
big steps: \absexpr{\emptyctx \types M \evalto v}, and evaluation contexts:
\absexpr{\ctxfill{C}{\hole}}.

\section{Outside of absint notation}
Symbols whose name is still free also work in plain math mode:
\(a \lub b \partof \top\), or with \verb|\absintuse|: \absintuse{widen}.

\section{Notation}
\absintlegend
\end{document}
